% !TEX program = xelatex
\documentclass[11pt,a4paper]{article}
\usepackage[margin=2.2cm,top=2cm]{geometry}
\usepackage{fontspec}
\usepackage{polyglossia}
\setdefaultlanguage{turkish}
\usepackage{amsmath,amssymb}
\usepackage{graphicx}
\usepackage{booktabs}
\usepackage{enumitem}
\usepackage{multicol}
\usepackage[most]{tcolorbox}
\usepackage{fancyhdr}
\usepackage{titlesec}
\usepackage{xcolor}
\usepackage{needspace}
\usepackage[hidelinks]{hyperref}

\definecolor{anamavi}{RGB}{0,45,98}
\definecolor{cozumyesil}{RGB}{0,110,50}
\definecolor{uyarikirmizi}{RGB}{180,20,20}

\titleformat{\section}{\large\bfseries\color{anamavi}}{\thesection.}{0.5em}{}
\titleformat{\subsection}{\normalsize\bfseries\color{anamavi}}{\thesubsection}{0.5em}{}

\pagestyle{fancy}\fancyhf{}
\lhead{\small MMAT105 Genel Matematik}
\rhead{\small Hafta 2 -- Ders Notu}
\cfoot{\small\thepage}
\renewcommand{\headrulewidth}{0.4pt}

\newtcolorbox{kural}[1][Kural]{colback=anamavi!5,colframe=anamavi,title=#1,fonttitle=\bfseries,breakable,left=4pt,right=4pt,top=2pt,bottom=2pt}
\newtcolorbox{dikkat}[1][Dikkat]{colback=uyarikirmizi!5,colframe=uyarikirmizi,title=#1,fonttitle=\bfseries,breakable,left=4pt,right=4pt,top=2pt,bottom=2pt}
\newcounter{ornek}
\newenvironment{ornek}{\refstepcounter{ornek}\par\medskip\noindent\textbf{\color{anamavi}Örnek \theornek.}\ }{\par}
\newenvironment{cozum}{\begin{tcolorbox}[colback=cozumyesil!6,colframe=cozumyesil,boxrule=0.5pt,left=4pt,right=4pt,top=2pt,bottom=2pt,breakable]\textbf{\color{cozumyesil}Çözüm.}\ }{\end{tcolorbox}}
\newcounter{alis}
\newenvironment{alistirma}{\refstepcounter{alis}\par\medskip\needspace{9\baselineskip}\noindent\textbf{\color{anamavi}Alıştırma \thealis.}\ }{\par}
\newcommand{\kendinizoncesi}{\needspace{12\baselineskip}}
\newcommand{\ab}[1]{\overline{#1}}
\setlength{\parindent}{0pt}
\setlength{\parskip}{3pt}

\begin{document}

% ---------------- Başlık ----------------
\begin{minipage}[c]{\textwidth}
{\LARGE\bfseries\color{anamavi} MMAT105 Genel Matematik}\\[4pt]
{\large Hafta 2 Ders Notu: Asal Sayılar, Asal Çarpanlara Ayırma,\\ Tam Bölenler ve Bölme}\\[6pt]
{\small 2026--2027 Güz}
\end{minipage}

\vspace{4pt}\noindent\textcolor{anamavi}{\rule{\textwidth}{1pt}}

\textbf{Bu notun kullanımı.} Her bölümde önce kısa bir kural kutusu, ardından çözümlü örnekler yer alır. Örnekleri çözümüne bakmadan önce kendiniz deneyin. Notun sonunda önce çözümlü alıştırmalar, ardından cevap anahtarlı ``Kendiniz Çözün'' alıştırmaları yer alır.

% =====================================================
\section{Asal Sayılar}

\begin{kural}[Tanım ve asallık testi]
\(1\)'den büyük olup \(1\) ve kendisinden başka pozitif böleni olmayan tam sayıya \textbf{asal sayı} denir. \(100\)'e kadar \(25\) asal sayı vardır:
\[2,3,5,7,11,13,17,19,23,29,31,37,41,43,47,53,59,61,67,71,73,79,83,89,97\]
\textbf{Test:} \(n>1\) sayısı, \(\sqrt{n}\)'den büyük olmayan asalların hiçbirine tam bölünmüyorsa asaldır.
\end{kural}

\begin{dikkat}
\(1\) asal değildir; \(0\) ve negatif sayılar asal değildir. \(2\) tek çift asaldır. İki asalın toplamı (veya farkı) tek ise bu asallardan biri \(2\)'dir.
\end{dikkat}

\begin{ornek}
\(60\) ile \(80\) arasında kaç asal sayı vardır? Bu asalların toplamı kaçtır?
\end{ornek}
\begin{cozum}
\(61,\ 67,\ 71,\ 73,\ 79\) \(\Rightarrow\) \(5\) tane. \quad Toplam: \(61+67+71+73+79=351\).
\end{cozum}

\begin{ornek}
\(307\) ve \(323\) sayıları asal mıdır?
\end{ornek}
\begin{cozum}
\(17^2=289\), \(18^2=324\): her iki sayı için \(2,3,5,7,11,13,17\) asallarını denemek yeter.\\
\(323=17\cdot 19\) \(\Rightarrow\) asal \textbf{değildir}. \quad \(307\) bu yedi asalın hiçbirine tam bölünmez \(\Rightarrow\) \textbf{asaldır}.
\end{cozum}

\begin{ornek}
\(a<b\) olmak üzere \(a\), \(b\) ve \(a+b\) sayılarının üçü de asaldır ve \(a+b<20\)'dir. Kaç farklı \((a,b)\) ikilisi vardır?
\end{ornek}
\begin{cozum}
\(a+b\) asal ve \(a+b>2\) olduğundan tektir; o halde \(a\) ile \(b\)'den biri çifttir: \(a=2\). \(2+b\) de asal olmalı:
\[(2,3)\to 5,\quad (2,5)\to 7,\quad (2,11)\to 13,\quad (2,17)\to 19\]
Cevap: \(4\) ikili. (\(b=7\Rightarrow 9\), \(b=13\Rightarrow 15\) asal değil.)
\end{cozum}

\begin{ornek}
\(a\) ve \(b\) asal sayılardır. \(2a+3b=23\) olduğuna göre \(a+b\) kaçtır?
\end{ornek}
\begin{cozum}
\(2a\) çift, \(23\) tek \(\Rightarrow\) \(3b\) tek \(\Rightarrow\) \(b\) tek asal. Deneyelim:
\(b=3\Rightarrow 2a=14\Rightarrow a=7\) (asal \checkmark); \ \(b=5\Rightarrow a=4\) (asal değil); \ \(b=7\Rightarrow a=1\) (asal değil). \(b\ge 11\) için \(3b>23\).
\[a+b=7+3=10\]
\end{cozum}

% =====================================================
\section{Aralarında Asal Sayılar}

\begin{kural}[Tanım]
\(1\)'den başka ortak pozitif böleni olmayan iki tam sayıya \textbf{aralarında asal} denir. Ardışık iki pozitif tam sayı ve farklı iki asal sayı her zaman aralarında asaldır.
\end{kural}

\begin{dikkat}
Aralarında asal sayıların kendilerinin asal olması gerekmez: \(8=2^3\) ve \(9=3^2\) aralarında asaldır.
\end{dikkat}

\begin{ornek}
\(1\)'den \(30\)'a kadar (\(1\) ve \(30\) dahil) olan tam sayılardan kaç tanesi \(15\) ile aralarında asaldır?
\end{ornek}
\begin{cozum}
\(15=3\cdot 5\): \(3\)'e veya \(5\)'e bölünen sayılar çıkarılır.\\
\(3\)'ün katları: \(10\) tane; \ \(5\)'in katları: \(6\) tane; \ ikisine birden bölünenler (\(15\), \(30\)) iki kez sayıldı: \(2\) tane.
\[30-(10+6-2)=30-14=16\]
\end{cozum}

% =====================================================
\section{Asal Çarpanlara Ayırma ve Bölen Sayısı}

\begin{kural}[Asal çarpanlar ve bölenler]
\(1\)'den büyük her tam sayı asal çarpanlarına tek türlü ayrılır (bölen merdiveni):
\[n=p_1^{\alpha_1}\cdot p_2^{\alpha_2}\cdots p_k^{\alpha_k}\qquad(p_i\ \text{farklı asallar},\ \alpha_i\ge 1)\]
\begin{itemize}[leftmargin=*,itemsep=1pt]
\item Pozitif bölen sayısı: \(d(n)=(\alpha_1+1)(\alpha_2+1)\cdots(\alpha_k+1)\); \quad tam bölen sayısı: \(2\cdot d(n)\)
\item Asal bölen sayısı: \(k\); \quad asal olmayan pozitif bölen sayısı: \(d(n)-k\)
\item Pozitif bölenler toplamı: \(\dfrac{p_1^{\alpha_1+1}-1}{p_1-1}\cdot\dfrac{p_2^{\alpha_2+1}-1}{p_2-1}\cdots\dfrac{p_k^{\alpha_k+1}-1}{p_k-1}\)
\end{itemize}
\end{kural}

\begin{dikkat}[Sık yapılan hatalar]
\(1\) bir asal çarpan değildir ama bir bölendir (asal olmayan bölen). Üsler toplanmaz; \((\alpha_i+1)\) çarpanları \textbf{çarpılır}. Tam bölenler negatif bölenleri de içerir.
\end{dikkat}

\begin{ornek}
\(1260\) sayısını asal çarpanlarına ayırınız. Farklı asal çarpanlarının toplamı ve pozitif bölen sayısı kaçtır?
\end{ornek}
\begin{cozum}
\begin{minipage}[t]{0.18\linewidth}
\begin{tabular}{r|l}
\(1260\) & \(2\)\\ \(630\) & \(2\)\\ \(315\) & \(3\)\\ \(105\) & \(3\)\\ \(35\) & \(5\)\\ \(7\) & \(7\)\\ \(1\) &
\end{tabular}
\end{minipage}\hfill
\begin{minipage}[t]{0.62\linewidth}
\[1260=2^2\cdot 3^2\cdot 5\cdot 7\]
Farklı asal çarpanlar toplamı: \(2+3+5+7=17\).
\[d(1260)=(2+1)(2+1)(1+1)(1+1)=36\]
\end{minipage}
\end{cozum}

\begin{ornek}
\(2^3\cdot 3^a\cdot 5\) sayısının \(24\) pozitif böleni vardır. \(a\) kaçtır?
\end{ornek}
\begin{cozum}
\[(3+1)(a+1)(1+1)=24\ \Rightarrow\ 8(a+1)=24\ \Rightarrow\ a=2\]
\end{cozum}

\begin{ornek}
\(10^n\) sayısının \(49\) pozitif böleni vardır. \(n\) kaçtır?
\end{ornek}
\begin{cozum}
\(10^n=2^n\cdot 5^n\) \(\Rightarrow\) \(d=(n+1)(n+1)=(n+1)^2=49\ \Rightarrow\ n+1=7\ \Rightarrow\ n=6\).
\end{cozum}

\begin{ornek}
\(540\) sayısının tam bölen sayısı ve asal olmayan pozitif bölen sayısı kaçtır?
\end{ornek}
\begin{cozum}
\(540=2^2\cdot 3^3\cdot 5\) \(\Rightarrow\) \(d=3\cdot 4\cdot 2=24\). \ Tam bölen: \(2\cdot 24=48\).\\
Asal bölenler \(2,3,5\) \(\Rightarrow\) asal olmayan pozitif bölen: \(24-3=21\).
\end{cozum}

\begin{ornek}
\(360\) sayısının pozitif bölenlerinden kaç tanesi \(6\)'nın katıdır?
\end{ornek}
\begin{cozum}
\(6\)'nın katı olan bir bölen \(6\cdot m\) biçimindedir ve \(6\cdot m\), \(360\)'ı bölüyorsa \(m\), \(360\div 6=60\)'ı böler.
\[60=2^2\cdot 3\cdot 5\ \Rightarrow\ d(60)=3\cdot 2\cdot 2=12\]
\end{cozum}

\begin{ornek}
\(36\) sayısının pozitif bölenlerinin toplamı kaçtır?
\end{ornek}
\begin{cozum}
\(36=2^2\cdot 3^2\): \quad \(\dfrac{2^3-1}{2-1}\cdot\dfrac{3^3-1}{3-1}=7\cdot 13=91\).\\
Kontrol: \(1+2+3+4+6+9+12+18+36=91\).
\end{cozum}

% =====================================================
\section{Tam Kare, Tam Küp ve \texorpdfstring{\(n!\)}{n!} İçindeki Asal Kuvvetleri}

\begin{kural}[Kural]
Asal çarpanlarına ayrılmış bir pozitif tam sayı \textbf{tam kare} ise bütün üsler çift, \textbf{tam küp} ise bütün üsler \(3\)'ün katıdır.\\
\(n!\) içindeki \(p\) asalının üssü: \(\left\lfloor\dfrac{n}{p}\right\rfloor+\left\lfloor\dfrac{n}{p^2}\right\rfloor+\left\lfloor\dfrac{n}{p^3}\right\rfloor+\cdots\) \quad (\(\lfloor x\rfloor\): \(x\)'in tam kısmı)
\end{kural}

\begin{ornek}
\(150\cdot n\) bir tam kare ve \(150\cdot m\) bir tam küp ise \(n\) ve \(m\) en küçük hangi pozitif tam sayılardır?
\end{ornek}
\begin{cozum}
\(150=2\cdot 3\cdot 5^2\).
Tam kare: \(2\) ve \(3\)'ün üsleri çifte tamamlanır: \(n=2\cdot 3=6\), \ \(150\cdot 6=900=30^2\).\\
Tam küp: üsler \(3\)'e tamamlanır: \(m=2^2\cdot 3^2\cdot 5=180\), \ \(150\cdot 180=27000=30^3\).
\end{cozum}

\begin{ornek}
\(25!\) sayısının asal çarpanlarına ayrılışında \(3\)'ün üssü kaçtır? \(25!\) sayısı \(6^k\) ile tam bölünüyorsa \(k\) en çok kaçtır?
\end{ornek}
\begin{cozum}
\(3\)'ün üssü: \(\lfloor 25/3\rfloor+\lfloor 25/9\rfloor=8+2=10\). \quad \(2\)'nin üssü: \(12+6+3+1=22\).\\
\(6^k=2^k\cdot 3^k\): az olan çarpan belirleyicidir \(\Rightarrow\) \(k\) en çok \(10\).
\end{cozum}

% =====================================================
\section{Bölme}

\begin{kural}[Bölme algoritması]
\(A,B,C,K\) doğal sayılar ve \(B\neq 0\) olmak üzere, \(A\)'nın \(B\) ile bölümünde bölüm \(C\), kalan \(K\) ise
\[A=B\cdot C+K,\qquad 0\le K<B.\]
\(K=0\) ise \(A\), \(B\)'ye tam bölünür. Kalanın alabileceği en büyük değer \(B-1\)'dir.
\end{kural}

\begin{dikkat}
Kalan bölenden küçük olmalıdır. Kalan \(K\) verilmişse bölen en az \(K+1\)'dir.
\end{dikkat}

\begin{ornek}
İki basamaklı bir doğal sayı \(9\) ile bölündüğünde bölüm ile kalanın toplamı \(12\) oluyor. Bu sayı en az ve en çok kaçtır?
\end{ornek}
\begin{cozum}
\(A=9C+K\), \ \(C+K=12\), \ \(0\le K\le 8\) \(\Rightarrow\) \(A=9C+(12-C)=8C+12\) ve \(C\ge 4\).\\
En az: \(C=4,\ K=8\Rightarrow A=44\). \quad En çok: \(8C+12\le 99\Rightarrow C\le 10\); \(C=10,\ K=2\Rightarrow A=92\).
\end{cozum}

\begin{ornek}
Bir doğal sayının \(12\) ile bölümünden kalan \(7\)'dir. Bu sayının \(4\) ile ve \(3\) ile bölümünden kalanlar kaçtır?
\end{ornek}
\begin{cozum}
\(A=12k+7\).
\[A=4(3k+1)+3\ \Rightarrow\ \text{kalan } 3,\qquad A=3(4k+2)+1\ \Rightarrow\ \text{kalan } 1\]
\end{cozum}

\begin{ornek}
\(A\)'nın \(4\) ile bölümünde bölüm \(B\), kalan \(3\); \ \(B\)'nin \(6\) ile bölümünde kalan \(5\)'tir. \(A\)'nın \(24\) ile bölümünden kalan kaçtır?
\end{ornek}
\begin{cozum}
\(B=6C+5\) \(\Rightarrow\) \(A=4(6C+5)+3=24C+23\). \ \(0\le 23<24\) \(\Rightarrow\) kalan \(23\).
\end{cozum}

\begin{ornek}
\(A\) doğal sayısı \(B\)'ye bölündüğünde bölüm \(5\), kalan \(7\)'dir. \(A+B=91\) olduğuna göre \(A\) kaçtır?
\end{ornek}
\begin{cozum}
\(A=5B+7\) \(\Rightarrow\) \(5B+7+B=91\Rightarrow 6B=84\Rightarrow B=14\), \ \(A=77\).\\
Kontrol: \(77=14\cdot 5+7\) ve \(0\le 7<14\) \checkmark.
\end{cozum}

\begin{ornek}
Bölenin \(8\) olduğu bir bölmede bölüm \(13\) ve kalan alabileceği en büyük değerdedir. Bölünen kaçtır?
\end{ornek}
\begin{cozum}
En büyük kalan \(8-1=7\). \quad \(A=8\cdot 13+7=111\).
\end{cozum}

% =====================================================
\section{Kalanlarla İşlem ve Birler Basamağı}

\begin{kural}[Kalanlarla işlem]
\(a\) ve \(b\)'nin \(m\) ile bölümünden kalanlar \(r_1\) ve \(r_2\) ise \(a+b\)'nin kalanı \(r_1+r_2\)'nin, \(a\cdot b\)'nin kalanı \(r_1\cdot r_2\)'nin \(m\) ile bölümünden kalandır. Bir doğal sayının birler basamağı, \(10\) ile bölümünden kalandır.
\end{kural}

\begin{ornek}
\(x\) doğal sayısının \(8\) ile bölümünden kalan \(5\)'tir. \(x^2+3x\) sayısının \(8\) ile bölümünden kalan kaçtır?
\end{ornek}
\begin{cozum}
\(x\) yerine kalanı yazılır: \(5^2+3\cdot 5=25+15=40=8\cdot 5\) \(\Rightarrow\) kalan \(0\).\\
Kontrol: \(x=13\): \(169+39=208=8\cdot 26\).
\end{cozum}

\begin{ornek}
\(13\cdot 17\cdot 19\cdot 23\) çarpımının ve \(7^{2026}\) sayısının birler basamağındaki rakamlar kaçtır?
\end{ornek}
\begin{cozum}
Birler basamakları: \(3\cdot 7=21\to 1\); \ \(1\cdot 9=9\); \ \(9\cdot 3=27\to 7\). \ Cevap: \(7\) \ (\(13\cdot 17\cdot 19\cdot 23=96577\)).\\
\(7\)'nin kuvvetlerinin birler basamakları \(7,9,3,1\) biçiminde \(4\)'lü tekrar eder. \(2026=4\cdot 506+2\) \(\Rightarrow\) ikinci terim: \(9\).
\end{cozum}

\begin{ornek}
Bugün cuma. \(250\) gün sonra hangi gün olur? \(8^{15}\) sayısının birler basamağındaki rakam kaçtır?
\end{ornek}
\begin{cozum}
\(250=7\cdot 35+5\): cumadan \(5\) gün sonra \(\Rightarrow\) \textbf{çarşamba}.\\
\(8\)'in kuvvetlerinin birler basamakları \(8,4,2,6\) biçiminde tekrar eder. \(15=4\cdot 3+3\) \(\Rightarrow\) üçüncü terim: \(2\).
\end{cozum}

\begin{ornek}
\(a\) doğal sayısının \(5\) ile bölümünden kalan \(3\)'tür. \(a^3\) sayısının \(5\) ile bölümünden kalan kaçtır?
\end{ornek}
\begin{cozum}
\(a^3=a\cdot a\cdot a\) \(\Rightarrow\) kalan, \(3\cdot 3\cdot 3=27\)'nin kalanıdır: \(27=5\cdot 5+2\) \(\Rightarrow\) kalan \(2\).\\
Kontrol: \(a=8\): \(512=5\cdot 102+2\).
\end{cozum}

\begin{ornek}\textit{(Karma örnek)}
\(p\) ve \(p^2+2\) sayılarının ikisi de asal ise \(p\) kaçtır?
\end{ornek}
\begin{cozum}
\(p=3\) ise \(p^2+2=11\) asaldır. \(p\neq 3\) ise \(p\), \(3\)'e bölünmez; kalanı \(1\) veya \(2\)'dir. Bu durumda \(p^2\)'nin \(3\) ile bölümünden kalan \(1\cdot 1=1\) ya da \(2\cdot 2=4\to 1\) olur. O halde \(p^2+2\)'nin kalanı \(1+2=3\to 0\): \(p^2+2\), \(3\)'e tam bölünür ve \(3\)'ten büyük olduğundan asal olamaz. Cevap: \(p=3\).
\end{cozum}

% =====================================================
\section{Özet: Formül Kartı}
\begin{center}
\renewcommand{\arraystretch}{1.6}
\begin{tabular}{@{}ll@{}}
\toprule
Konu & Formül / Kural\\\midrule
Asal sayı & \(p>1\), pozitif bölenleri yalnız \(1\) ve \(p\); \ \(1\) asal değil, \(2\) tek çift asal\\
Asallık testi & \(\sqrt{n}\)'e kadar olan asalların hiçbiri bölmüyorsa \(n\) asal\\
Aralarında asal & ortak pozitif bölen yalnız \(1\); ardışık sayılar hep aralarında asal\\
Asal çarpanlar & \(n=p_1^{\alpha_1}\cdots p_k^{\alpha_k}\) (tek türlü)\\
Bölen sayısı & \(d(n)=(\alpha_1+1)\cdots(\alpha_k+1)\); \ tam bölen \(2\,d(n)\); \ asal olmayan \(d(n)-k\)\\
Bölenler toplamı & \(\prod \dfrac{p_i^{\alpha_i+1}-1}{p_i-1}\)\\
Tam kare / küp & üsler çift / üsler \(3\)'ün katı\\
\(n!\) içinde \(p\) & \(\lfloor n/p\rfloor+\lfloor n/p^2\rfloor+\cdots\)\\
Bölme & \(A=B\cdot C+K\), \ \(0\le K<B\); \ en büyük kalan \(B-1\)\\
Kalanlar & toplam/çarpımın kalanı \(=\) kalanların toplamı/çarpımının kalanı\\
Birler basamağı & \(10\) ile bölümden kalan; \(2,3,7,8\) kuvvetlerinde \(4\)'lü tekrar\\\bottomrule
\end{tabular}
\end{center}

% =====================================================
\newpage
% Hafta 2 -- alıştırma gövdesi (ders notu ve ayrı alıştırma kağıdı tarafından \input edilir)
\section{Çözümlü Alıştırmalar}
Her soruyu önce kendiniz deneyin, sonra çözümü okuyun.

\begin{alistirma}
\(p\) ve \(p+9\) sayılarının ikisi de asal ise \(p+(p+9)\) toplamı kaçtır?
\end{alistirma}
\begin{cozum}
\(p\) tek bir asal olsaydı \(p+9\) çift ve \(2\)'den büyük olurdu, yani asal olamazdı. O halde \(p\) çift asaldır: \(p=2\).\\
\(p+9=11\) asal \(\checkmark\). \quad Toplam: \(2+11=13\).
\end{cozum}

\begin{alistirma}
\(1\)'den \(24\)'e kadar (\(1\) ve \(24\) dahil) olan tam sayılardan kaç tanesi \(10\) ile aralarında asaldır?
\end{alistirma}
\begin{cozum}
\(10=2\cdot 5\): sayı ne \(2\)'ye ne \(5\)'e bölünmeli.\\
\(2\)'nin katları: \(12\) tane; \ \(5\)'in katları: \(5,10,15,20\) \(\to\) \(4\) tane; \ ikisinin de katı (\(10,20\)): \(2\) tane.\\
\(2\)'ye veya \(5\)'e bölünenler: \(12+4-2=14\). \quad Aralarında asal olanlar: \(24-14=10\).\\
(Liste: \(1,3,7,9,11,13,17,19,21,23\).)
\end{cozum}

\begin{alistirma}
\(120\) sayısının asal olmayan kaç pozitif böleni vardır?
\end{alistirma}
\begin{cozum}
\(120=2^3\cdot 3\cdot 5\) \(\Rightarrow\) pozitif bölen sayısı \((3+1)(1+1)(1+1)=16\).\\
Asal bölenler: \(2,3,5\) \(\to\) \(3\) tane. \quad Asal olmayanlar: \(16-3=13\) \ (\(1\) de asal değildir, bu sayıya dahildir).
\end{cozum}

\begin{alistirma}
\(90\) sayısının pozitif bölenlerinden kaç tanesi çift sayıdır?
\end{alistirma}
\begin{cozum}
\(90=2\cdot 3^2\cdot 5\) \(\Rightarrow\) bölen sayısı \(2\cdot 3\cdot 2=12\).\\
Tek bölenlerde \(2\) çarpanı bulunmaz: \(3^2\cdot 5\)'in bölen sayısı \(3\cdot 2=6\).\\
Çift bölenler: \(12-6=6\) \ (\(2,6,10,18,30,90\)).
\end{cozum}

\begin{alistirma}
\(108\cdot n\) bir tam küp ise \(n\) en küçük hangi pozitif tam sayıdır?
\end{alistirma}
\begin{cozum}
\(108=2^2\cdot 3^3\). Tam küpte bütün üsler \(3\)'ün katı olmalıdır.\\
\(2^2\to 2^3\) için bir \(2\) gerekir; \(3^3\) zaten uygundur. \quad \(n=2\), \ \(108\cdot 2=216=6^3\).
\end{cozum}

\begin{alistirma}
Bir doğal sayı \(15\) ile bölündüğünde bölüm \(7\) oluyor. Bu sayı en az ve en çok kaçtır?
\end{alistirma}
\begin{cozum}
\(A=15\cdot 7+K=105+K\), \quad \(0\le K\le 14\).\\
En az: \(K=0\Rightarrow A=105\). \quad En çok: \(K=14\Rightarrow A=119\).
\end{cozum}

\begin{alistirma}
\(A\)'nın \(7\) ile bölümünde bölüm \(B\), kalan \(4\); \ \(B\)'nin \(3\) ile bölümünde bölüm \(C\), kalan \(2\)'dir. \(A\)'nın \(21\) ile bölümünden kalan kaçtır?
\end{alistirma}
\begin{cozum}
\(A=7B+4\) \ ve \ \(B=3C+2\). \ Yerine yazalım:
\[A=7(3C+2)+4=21C+14+4=21C+18\]
\(18<21\) olduğundan kalan \(18\)'dir. \ (Kontrol: \(C=0\Rightarrow B=2,\ A=18\).)
\end{cozum}

\begin{alistirma}
Pozitif bir tam sayının \(9\) ile bölümünde bölüm, kalanın \(2\) katıdır. Bu koşulu sağlayan en büyük sayı kaçtır?
\end{alistirma}
\begin{cozum}
Kalan \(K\), bölüm \(2K\): \ \(A=9\cdot 2K+K=19K\).\\
\(0\le K\le 8\); sayı pozitif olduğundan \(K\ge 1\). En büyük \(A\) için \(K=8\):
\[A=19\cdot 8=152\qquad(\text{kontrol: } 152=9\cdot 16+8,\ 16=2\cdot 8)\]
\end{cozum}

\kendinizoncesi
\section{Alıştırmalar (Kendiniz Çözün)}

\begin{enumerate}[label=\textbf{\arabic*.},leftmargin=*,itemsep=6pt]
\item \(30\) ile \(60\) arasında kaç asal sayı vardır?
\item \(187\) ve \(191\) sayıları asal mıdır?
\item \(a\) ve \(b\) asal sayılardır. \(a+b=31\) olduğuna göre \(a\cdot b\) kaçtır?
\item \(252\) sayısını asal çarpanlarına ayırınız. Kaç pozitif böleni vardır?
\item \(400\) sayısının kaç tam böleni vardır?
\item \(2^a\cdot 5^3\) sayısının \(20\) pozitif böleni vardır. \(a\) kaçtır?
\item \(18\) sayısının pozitif bölenlerinin toplamı kaçtır?
\item \(360\) sayısının kaç tane tek pozitif böleni vardır?
\item \(30!\) sayısının asal çarpanlarına ayrılışında \(5\)'in üssü kaçtır?
\item Bir doğal sayının \(6\) ile bölümünden kalan \(4\)'tür. Bu sayının \(5\) katının \(6\) ile bölümünden kalan kaçtır?
\item \(a\)'nın \(11\) ile bölümünden kalan \(6\), \(b\)'ninki \(9\)'dur. \(a+b\) ve \(a\cdot b\)'nin \(11\) ile bölümünden kalanlar kaçtır?
\item \(3^{50}\) sayısının birler basamağındaki rakam kaçtır?
\end{enumerate}

\subsection*{Cevap Anahtarı}
{\small
\begin{multicols}{2}
\begin{enumerate}[label=\textbf{\arabic*.},leftmargin=*,itemsep=1pt]
\item \(31,37,41,43,47,53,59\): \(7\) tane.
\item \(187=11\cdot 17\) asal değil; \(191\) asal.
\item \(2+29=31\Rightarrow a\cdot b=58\).
\item \(252=2^2\cdot 3^2\cdot 7\); \(18\) bölen.
\item \(400=2^4\cdot 5^2\); \(15\) pozitif, \(30\) tam bölen.
\item \((a+1)\cdot 4=20\Rightarrow a=4\).
\item \(1+2+3+6+9+18=39\).
\item \(360=2^3\cdot 3^2\cdot 5\); tek bölen \(3\cdot 2=6\).
\item \(\lfloor 30/5\rfloor+\lfloor 30/25\rfloor=6+1=7\).
\item \(5\cdot 4=20=18+2\Rightarrow\) kalan \(2\).
\item \(a+b\to 4\); \ \(a\cdot b\to 10\).
\item \(3,9,7,1\) tekrarı; \(50=4\cdot 12+2\Rightarrow 9\).
\end{enumerate}
\end{multicols}}


\vspace{6pt}
\noindent\textcolor{anamavi}{\rule{\textwidth}{0.5pt}}\\
{\small Gelecek hafta: Bölünebilme kuralları, EBOB ve EKOK.}

\end{document}
